\section{未来研究方向与个人思考}
\label{sec:future}

\subsection{面向主导性能瓶颈的自适应放置框架}

本文设想一个简洁的自适应放置框架。每个决策窗口记录跨节点 All-to-All 暴露等待、各 expert 负载、sample routing locality、热点持续时间和可用显存，先判断主要约束是网络路径还是持续性专家计算，再评估 sample placement 或 expert replication；当专家权重较小且可复用时，把 Janus 类 fetch 作为改变通信对象的备选。通信调度和 pipeline 作为共同执行层，不作为互斥的第三类放置。

该框架不增加复杂控制器，而把观察、诊断、选择分成可检查步骤：先产生统一统计，再判断网络路径、专家计算或调整控制是否占据关键路径，最后只在成本条件满足时改变对象位置。这样既保留固定布局基线，也让样本重排和专家复制对应不同瓶颈。

\subsection{考虑动态调整成本的策略选择}

该框架的选择规则是：只有当目标窗口预计减少的关键路径暴露时间，大于无法被计算隐藏的求解、状态同步、数据搬运和切换成本，且局部性或热点足以摊销调整时才重配置；否则保持当前布局。该规则作为后续实验的选择假设。

为避免边界附近反复切换，可设置触发与回退阈值：连续多个窗口满足条件才重排或扩展副本，不满足条件才回退。阈值只稳定决策，不替代真实暴露时间和同步成本的测量。作为设计原则，路由快速漂移时可保持布局并调度通信；热点持续且显存充足时再复制；局部性强时评估样本重排。

\subsection{实验设计与验证方法}

验证分为两个层次。层 A 是控制变量测试：固定同一 MoE 模型、Top-k、硬件拓扑和路由 trace，设置固定专家并行、样本放置和动态专家复制基准；通信调度、token 接纳规则及其他通信设置保持一致，观察通信暴露、step time、专家热点、调整开销和显存变化。该层用于隔离放置动作，不能推出长期收敛结论。

实验交叉变化两条压力轴：高或低路由局部性，以及稳定或快速变化的专家热点，同时改变调整频率或触发阈值。

层 B 是端到端训练验证：允许 Router、模型状态和 expert capacity 随训练更新，考察 token 丢弃、训练正确性以及达到相同目标 loss 或质量的 time-to-convergence。若讨论精度、收敛或 token 丢弃，必须在相同目标质量下记录该指标，并记录动态副本的 per-step 与同步开销。层 B 不能用固定 trace 重放代替，因为它不包含路由状态、token 接纳和收敛反馈。

两层实验均记录 step time 及 P95、All-to-All 暴露等待、专家负载、跨节点字节、重配置/同步时间或字节、峰值显存和样本/序列语义正确性。可证伪判据包括：样本重排若减少跨节点通信却未改善端到端时间，须分解暴露、重排求解/同步、计算和尾部等待；复制若缓解热点却未加速，须检查同步、权重更新和重配置是否可见。高局部性且负载均衡用于检验样本重排，低局部性且热点稳定用于检验复制，快速变化用于检验阈值和保持现状。

为检查上述选择条件，本文另用 Python 运行合成 Trace 机制仿真：2 个 Node、4 张逻辑 GPU、8 个 Expert、12 个窗口、3 个 seed，比较 Keep、sample placement、expert replication 和 Cost-aware。结果仅在相对成本模型中成立：高局部性场景下样本放置降低跨节点 Token，稳定热点场景下复制降低峰值 GPU Token 负载但增加同步成本；快速热点漂移场景中，Cost-aware 依赖上一窗口信息，最终与样本放置一致，3 个 seed 的平均成本代理为 1011.68，高于 Keep 的 847.42（相对模型单位，仅作机制检查）。场景 A 的 3 个 seed 样本—专家计数相同，仅 Token 顺序不同，不是独立负载复现。场景 C 损失主要来自过时路由预测造成跨节点 Token 增多，而非复制/同步或单纯重配置费用；这表明简化策略缺少漂移检测与回退，支持稳定性判据，但未验证理想方案。这里的数值不是 GPU Step Time、真实网络测量或收敛收益。

